\section{Đường gradient ngắn lại bao nhiêu}
\label{sec:ch06-duong-gradient}

\index{attention!đường gradient thẳng}%
Phương trình~\eqref{eq:ch06-tong-co-trong-so} nhìn như một mẹo kỹ thuật. Nó
không phải. Nó đổi hình dạng đồ thị tính toán, và chỗ đổi hiện ra rõ nhất khi
lấy đạo hàm. Mục này dẫn ra chỗ ấy rồi đo nó.

Câu hỏi cụ thể: mất mát ở một bước đích \(i\) gửi được bao nhiêu tín hiệu về
tới annotation của vị trí nguồn \(j\)? Viết \(L_i\) cho mất mát ở riêng bước
\(i\).

\subsection*{Trong mô hình chương trước}

\index{attention!đường gradient của vector cố định}%
Ở chương~\ref{ch:encoder-decoder}, decoder chỉ nhìn thấy câu nguồn qua đúng một
vật, là trạng thái cuối cùng của encoder. Trong ký hiệu của chương này,
\(\vect{v} = \vect{h}_{T}\). Nên mọi đường từ \(L_i\) về \(\vect{h}_j\) đều
phải đi qua \(\vect{h}_{T}\):

\begin{equation}
  \pd{L_i}{\vect{h}_j}
  = \pd{L_i}{\vect{h}_{T}} \prod_{t = j+1}^{T} \jac{\vect{h}_t}{\vect{h}_{t-1}}.
  \label{eq:ch06-duong-cu}
\end{equation}

Cái tích ở vế phải là đúng vật mà chương~\ref{ch:phan-tich} đã dành cả chương
để phân tích. Nó có \(T - j\) thừa số. Nếu \tn{giá trị kỳ dị}{singular value}
lớn nhất của Jacobi hồi quy nhỏ hơn 1 thì chuẩn của tích giảm theo cấp số nhân
theo \(T - j\), và đó là \tn{điều kiện đủ}{sufficient condition} cho
\tn{triệt tiêu gradient}{vanishing gradient} mà
mục~\ref{sec:ch03-tich-jacobian} chứng minh.

Đọc \eqref{eq:ch06-duong-cu} theo nghĩa thường: từ nguồn càng nằm xa cuối câu
thì tín hiệu học của nó càng phải đi qua nhiều thừa số. Từ đầu câu chịu nhiều
nhất.

\subsection*{Trong mô hình chương này}

\index{attention!đạo hàm qua tổng có trọng số}%
Bây giờ \(\vect{h}_j\) vào thẳng \(\vect{c}_i\) qua
\eqref{eq:ch06-tong-co-trong-so}. Lấy đạo hàm tổng ấy theo một số hạng của nó:

\begin{equation}
  \jac{\vect{c}_i}{\vect{h}_j}
  = \alpha_{ij} \mat{I}
  + \sum_{k=1}^{T} \vect{h}_k \, \pd{\alpha_{ik}}{\vect{h}_j}\transpose .
  \label{eq:ch06-duong-moi}
\end{equation}

Số hạng thứ hai có mặt vì \(\vect{h}_j\) cũng tham gia chấm điểm, nên nó tác
động lên cả các trọng số. Số hạng \emph{thứ nhất} mới là chỗ cần nhìn:
\(\alpha_{ij} \mat{I}\). Một số vô hướng nhân ma trận đơn vị. Không có tích
Jacobi nào, không có thừa số nào phụ thuộc \(T - j\), không có gì để nhân dồn.

Đặt hai vế cạnh nhau thì phát biểu được gọn:

\begin{quote}
Ở chương trước, đường ngắn nhất từ bước đích \(i\) về annotation \(j\) dài
\(T - j\) bước. Ở đây nó dài đúng một bước, với mọi \(j\).
\end{quote}

\index{attention!chỗ dẫn xuất này không nói}%
Ba chỗ dẫn xuất này \emph{không} nói, và bỏ qua chúng là cách đọc quá tay.

Thứ nhất, đường dài cũ vẫn còn nguyên. Attention không cắt đường nào cả, nó
\emph{thêm} một đường ngắn vào bên cạnh. Gradient vẫn chảy dọc encoder như cũ;
chỉ là bây giờ nó có thêm lối khác để đi.

Thứ hai, đường ngắn ấy dẫn tới \emph{annotation}, không dẫn tới từ. Bản thân
\(\vect{h}_j\) do một mạng hồi quy tính ra, nên nó vẫn phụ thuộc vào cả đoạn
câu trước nó. Cái đường một bước rút ngắn là đoạn từ bước đích tới annotation,
không phải toàn bộ đoạn từ bước đích tới từ.

Thứ ba, hệ số \(\alpha_{ij}\) nằm ngay trước \(\mat{I}\). Nếu mô hình đặt
\(\alpha_{ij}\) gần 0 thì tín hiệu về \(\vect{h}_j\) qua đường ngắn cũng gần 0.
Khác biệt với chương trước là ở chỗ con số ấy do mô hình chọn theo từng bước,
chứ không do khoảng cách quyết định.

\subsection*{Đo}

\index{attention!đo tầm với của gradient}%
Dẫn xuất nói đường ngắn tồn tại. Nó không nói tín hiệu thật sự đi tới được bao
nhiêu, vì trong \eqref{eq:ch06-duong-cu} còn có hằng số, có phương của vector,
có đủ thứ mà giấy không nói nổi. Nên đo.

Cách đo: dựng một mô hình mỗi loại, cùng seed nên embedding xuất phát giống hệt
nhau, lấy mất mát ở riêng một bước đích, \tn{lan truyền ngược}{backpropagation},
rồi đọc chuẩn của
gradient tới embedding của từng từ nguồn. Chia mỗi cột cho phần tử lớn nhất của
chính nó, vì hai kiến trúc có thang khác nhau và cái tôi hỏi là \emph{hình
dạng}. Một cột phẳng nghĩa là mọi từ nguồn nhận được cùng cỡ tín hiệu; một cột
dốc nghĩa là đầu câu bị bỏ lại.

\index{attention!đo tại lúc khởi tạo}%
Đo \textbf{tại lúc khởi tạo, chưa huấn luyện gì cả}. Đây là lựa chọn có chủ ý
và nó là lựa chọn chương~\ref{ch:lstm} đã làm một lần rồi cho đạo hàm qua
vòng quay lỗi. Thứ đang kiểm là một tính chất của đồ thị tính toán, tức đường
nào tồn tại và chúng làm gì với gradient; còn huấn luyện là thứ xảy ra
\emph{nhờ} các đường ấy. Đo sau khi huấn luyện thì con số trả lời một câu hỏi
khác, là corpus này đã dạy được mô hình này cái gì.

Và nhờ chi tiết \(\vect{v}_a = 0\) ở cuối mục trước, điểm xuất phát sạch: mọi
\(\alpha_{ij}\) bằng \(1/T\). Mô hình attention chưa hề ưu tiên vị trí nào.

\begin{measured}{tầm với của gradient theo độ dài câu, chưa huấn luyện}
\begin{minted}{text}
source_len  fixed min/max  attention min/max  ratio
5           0.237278       0.853556           3.6x
9           0.113478       0.626161           5.5x
14          0.028660       0.554965           19.4x
18          0.012126       0.604488           49.9x
21          0.005483       0.586938           107.0x
\end{minted}

\(d_h = 32\), mất mát lấy ở bước đích 5, mỗi dòng một câu trong tập test.
\code{min/max} là chuẩn gradient nhỏ nhất chia cho lớn nhất trên các vị trí
nguồn: 1.0 là phẳng tuyệt đối.

Đo bằng \code{experiments/ch06\_gradient.py} tại tag \code{ch06}.
\end{measured}

Cột \code{fixed} là hình dạng của \eqref{eq:ch06-duong-cu} khi vẽ ra. Câu 5
token còn giữ được 0.237278, câu 21 token còn 0.005483. Hơn bốn mươi ba lần
suy giảm, và thứ duy nhất thay đổi giữa hai dòng là câu dài thêm mười sáu
token.

Cột \code{attention} không có xu hướng nào. Nó đi 0.853556, 0.626161, 0.554965,
0.604488, 0.586938 và cứ đứng quanh đó. Bốn dòng cuối nằm gọn giữa 0.554965 và
0.626161, và chúng lên xuống không theo chiều nào của độ dài câu.

\index{attention!vì sao tỷ số mới là số đáng đọc}%
Cột \code{ratio} là chỗ tôi muốn bạn dừng lại. Không phải vì attention phẳng
hơn, chuyện đó thì \eqref{eq:ch06-duong-moi} đã nói trước rồi. Mà vì khoảng
cách giữa hai kiến trúc \emph{lớn dần theo độ dài câu}: 3.6 lần ở câu ngắn,
107.0 lần ở câu dài nhất trong tập test.

Đó là phát biểu sắc hơn \enquote{attention giúp gradient đi xa hơn}. Nó nói
rằng cái giá của vector cố định không phải một hằng số phải trả, nó là một hàm
tăng theo độ dài câu. Đặt cạnh hình 4 của bài Cho 2014b ở mục~\ref{sec:ch06-ai-do},
hai thứ khớp nhau: một đường BLEU tụt theo độ dài câu, và một cơ chế mà tầm với
của nó co lại theo đúng độ dài ấy.

\subsection*{Nhìn từng từ một}

\index{attention!hình dạng theo vị trí từ nguồn}%
Bảng trên nén mỗi câu thành một số. Mở một câu ra thì thấy hình dạng.

\begin{measured}{gradient tới từng từ nguồn, câu test dài nhất}
\begin{minted}{text}
j   source word   fixed     attention
1   a             0.0055    0.6493
2   big           0.0069    0.6951
3   cat           0.0069    0.7488
4   carries       0.0130    0.9483
5   a             0.0134    0.6718
6   white         0.0162    0.7804
7   book          0.0229    0.7809
8   in            0.0336    0.7891
9   the           0.0481    0.7796
10  garden        0.0471    0.7056
11  and           0.0734    0.8282
12  a             0.0803    0.7465
13  small         0.1111    0.7344
14  dog           0.1781    0.8336
15  carries       0.2186    1.0000
16  a             0.2391    0.6914
17  big           0.3990    0.7692
18  cat           0.3895    0.8106
19  at            0.6097    0.7753
20  the           0.7639    0.6209
21  window        1.0000    0.5869
\end{minted}

Mỗi cột chia cho phần tử lớn nhất của chính nó. Cùng lần chạy với bảng trên.

Đo bằng \code{experiments/ch06\_gradient.py} tại tag \code{ch06}.
\end{measured}

Cột \code{fixed} tăng gần như đơn điệu từ trên xuống dưới và chỉ đạt 1.0000 ở
token cuối cùng, đúng chỗ encoder dừng lại. Từ \code{window} ở vị trí 21 nhận
được nhiều hơn từ \code{a} ở vị trí 1 khoảng một trăm tám mươi lần. Cả hai đều
là từ của cùng một câu, và câu ấy chỉ có hai mệnh đề.

Cột \code{attention} không xếp theo vị trí. Giá trị lớn nhất rơi vào vị trí 15,
tức \code{carries}, nằm giữa câu. Không có lý do sâu xa nào ở đó: mô hình chưa
huấn luyện, các trọng số bằng nhau hết, và cái còn lại là nhiễu của khởi tạo.
Điều đáng nói chính là chuyện ấy. Ở cột bên trái, thứ tự các con số do
\emph{vị trí} quyết định. Ở cột bên phải, vị trí không quyết định gì.

\index{attention!sau khi huấn luyện thì khác}%
Một chú ý cuối cùng, để không ai mang bảng này đi quá xa. Sau khi huấn luyện,
cột \code{attention} sẽ \emph{không} còn phẳng. Mô hình học được cách dồn
\(\alpha_{ij}\) vào một hai vị trí, nên gradient cũng dồn theo, và
mục~\ref{sec:ch06-mo-ma-tran} in ra đúng những trọng số dồn ấy. Chỗ khác nhau
nằm ở nguyên nhân: ở đây một cột dốc là mô hình đang chọn, còn ở cột
\code{fixed} một cột dốc là khoảng cách đang áp đặt. Hai thứ trông giống nhau
trên biểu đồ và không giống nhau chút nào.
